\chapter*{ABSTRAK}
\addcontentsline{toc}{chapter}{Abstrak} 

Perkembangan aplikasi web yang pesat meningkatkan risiko munculnya celah keamanan, seperti \textit{SQL Injection}, \textit{Cross-Site Scripting} (XSS), dan \textit{basic fuzzing}. Metode \textit{automated web penetration testing} konvensional yang masih berbasis aturan statis (\textit{rule-based}) memiliki keterbatasan dalam aspek adaptivitas dan efisiensi, terutama ketika menghadapi lingkungan web yang dinamis. Penelitian ini bertujuan untuk merancang dan mengimplementasikan metode \textit{Reinforcement Learning} dengan algoritma \textit{Q-Learning} dalam sistem \textit{automated web penetration testing}. Sistem dirancang agar agen dapat memilih strategi pengujian berdasarkan respons dari aplikasi target melalui mekanisme \textit{state}, \textit{action}, \textit{reward}, dan pembaruan \textit{Q-Table}. Pengujian dilakukan pada \textit{OWASP Juice Shop} dan dibandingkan dengan baseline menggunakan \textit{OWASP ZAP}. Evaluasi dilakukan berdasarkan jumlah kerentanan yang ditemukan, jumlah \textit{request}, waktu pengujian, \textit{precision}, dan \textit{cumulative reward}. Hasil yang diharapkan adalah sistem pengujian keamanan web otomatis yang lebih adaptif dan efisien dibandingkan dengan metode konvensional.

\textbf{Kata kunci:} \textit{penetration testing, reinforcement learning, Q-learning, web security, automated testing, OWASP Juice Shop.}
\clearpage


\chapter*{ABSTRACT}
\addcontentsline{toc}{chapter}{Abstract} 

The rapid development of web applications increases the risk of security vulnerabilities, such as SQL Injection, Cross-Site Scripting (XSS), and basic fuzzing. Conventional automated web penetration testing methods that rely on static rules are limited in adaptability and efficiency, especially in dynamic web environments. This study aims to design and implement a Reinforcement Learning method based on the Q-Learning algorithm for an automated web penetration testing system. The system is designed to enable the agent to select testing strategies based on responses from the target application via state, action, reward, and Q-Table updates. Testing is conducted on OWASP Juice Shop and compared with a baseline using OWASP ZAP. The evaluation is based on the number of detected vulnerabilities, the number of requests, testing time, precision, and cumulative reward. The expected result is an automated web security testing system that is more adaptive and efficient than conventional methods.

\textit{\textbf{Keywords:} penetration testing, reinforcement learning, Q-learning, web security, automated testing, OWASP Juice Shop.}
\clearpage